%! Polynomial Functions and End Behavior — Unit 2, Lesson 2.5 — Student Packet Cover Sheet
\documentclass[10pt]{article}
\usepackage{precalculus-article}
\usepackage{precalculus-boxes}
\usepackage{ltablex}
\keepXColumns

\begin{document}

% Full-bleed plum banner
\begin{tikzpicture}[remember picture, overlay]
  \fill[plum] (current page.north west)
    rectangle ([yshift=-0.9in]current page.north east);
\end{tikzpicture}
\vspace{-0.8in}

\begin{center}
  {\color{white}\textbf{\LARGE Precalculus}}\\[4pt]
  {\color{lilacmid}\large\textbf{Unit 2: Polynomial Functions}}\\[6pt]
  {\color{white}\textbf{\large Lesson 2.5 \quad Polynomial Functions and End Behavior}}
\end{center}

\vspace{0.22in}
\namedateperiod
\vspace{0.18in}

\begin{learningtargetbox}
\small
\begin{enumerate}[leftmargin=1.5em, itemsep=5pt]
  \item I can say in words what it means for the output values of a polynomial
        to increase or decrease without bound, and I can write that description
        in limit notation such as $\lim_{x \to \infty} p(x) = \infty$.
  \item I can show with numbers that the leading term of a polynomial dominates
        every lower-degree term once the inputs are large, and I can use that
        fact to explain why the leading term alone decides end behavior.
  \item I can use the degree and the sign of the leading coefficient to place a
        polynomial in one of the four end-behavior cases and write both limit
        statements for it.
  \item I can find the leading term of a polynomial written in standard form or
        in factored form, and I can read end behavior off a graph.
\end{enumerate}
\end{learningtargetbox}

\vspace{0.14in}

\begin{tocbox}
\small
\renewcommand{\arraystretch}{1.5}
\begin{tabularx}{\linewidth}{@{} c l X r @{}}
\rowcolor{lilac}
\textbf{\#} & \textbf{Component} & \textbf{Description} & \textbf{Score} \\
\midrule
1 & Warm-Up        & Spiral review of prerequisite skills                        & \blank{1.2cm} \\
2 & Guided Notes   & Limit notation, dominance, and the four end-behavior cases  & \blank{1.2cm} \\
3 & Group Activity & The End Behavior Desk: predicting both far ends             & \blank{1.2cm} \\
4 & Exit Ticket    & Independent check on limit notation and leading terms       & \blank{1.2cm} \\
5 & Homework       & Practice and extension                                      & \blank{1.2cm} \\
\midrule
  & \multicolumn{2}{r}{\textbf{Total}} & \blank{1.2cm} \\
\end{tabularx}
\end{tocbox}

\vspace{0.14in}

\begin{remindbox}
\small
Last lesson was about the middle of the graph --- the place where the zeros
live. Today is about the two far ends, and they turn out to be decided by a
single term. Once $|x|$ is large enough, every lower-degree term is a rounding
error beside the leading one, so the \textbf{degree} and the \textbf{sign of the
leading coefficient} are the only two facts you need.
\end{remindbox}

\end{document}
